%==============================================================================
%  ESTILO DE PÁGINA, TÍTULOS Y ENTORNOS DE LA OBRA
%==============================================================================

% --- Caja de texto --------------------------------------------------------
\setstocksize{240mm}{170mm}
\settrimmedsize{240mm}{170mm}{*}
\settypeblocksize{188mm}{116mm}{*}
\setlrmargins{*}{*}{1.35}      % margen exterior mayor (encuadernación)
\setulmargins{*}{*}{1.15}
\setheadfoot{\onelineskip}{2\onelineskip}
\setheaderspaces{*}{*}{1.2}
\checkandfixthelayout

\setSingleSpace{1.06}
\SingleSpacing

% --- Estilo de capítulo ---------------------------------------------------
\newcommand{\ReglaCapitulo}{%
  {\color{ColorBorde}\rule{\textwidth}{0.6pt}}}

\makechapterstyle{concursal}{%
  \setlength{\beforechapskip}{28pt}
  \setlength{\midchapskip}{10pt}
  \setlength{\afterchapskip}{34pt}
  \renewcommand{\chapnamefont}{\normalfont\sffamily\small\bfseries\color{ColorGris}}
  \renewcommand{\chapnumfont}{\normalfont\sffamily\small\bfseries\color{ColorGris}}
  \renewcommand{\chaptitlefont}{\normalfont\sffamily\LARGE\bfseries\color{ColorTinta}\raggedright}
  \renewcommand{\printchaptername}{\chapnamefont\MakeUppercase{\chaptername}}
  \renewcommand{\chapternamenum}{\space}
  \renewcommand{\afterchapternum}{\par\vskip\midchapskip}
  \renewcommand{\printchaptertitle}[1]{\chaptitlefont ##1\par\vskip6pt\ReglaCapitulo}
}
\chapterstyle{concursal}

% --- Estilo de secciones --------------------------------------------------
\setsecheadstyle{\normalfont\sffamily\large\bfseries\color{ColorTinta}\raggedright}
\setsubsecheadstyle{\normalfont\sffamily\normalsize\bfseries\color{ColorTinta}\raggedright}
\setsubsubsecheadstyle{\normalfont\sffamily\normalsize\itshape\color{ColorGris}\raggedright}
\setparaheadstyle{\normalfont\normalsize\bfseries}

\setsecnumdepth{subsection}
\settocdepth{subsection}

% --- Encabezados y pies ---------------------------------------------------
\makepagestyle{concursal}
\makeevenhead{concursal}{\sffamily\small\color{ColorGris}\thepage}{}%
  {\sffamily\small\color{ColorGris}\leftmark}
\makeoddhead{concursal}{\sffamily\small\color{ColorGris}\rightmark}{}%
  {\sffamily\small\color{ColorGris}\thepage}
\makeheadrule{concursal}{\textwidth}{0.4pt}
\makeevenfoot{concursal}{}{}{}
\makeoddfoot{concursal}{}{}{}
\pagestyle{concursal}

\renewcommand{\chaptermark}[1]{\markboth{\MakeUppercase{#1}}{}}
\renewcommand{\sectionmark}[1]{\markright{\thesection\quad #1}}

% --- Notas al pie ---------------------------------------------------------
\setlength{\footmarkwidth}{1.2em}
\setlength{\footmarksep}{-\footmarkwidth}
\footmarkstyle{\textsuperscript{#1}\hspace{0.35em}}
\renewcommand{\foottextfont}{\footnotesize}

% --- Párrafos -------------------------------------------------------------
\setlength{\parindent}{1.2em}
\nonzeroparskip
\setlength{\parskip}{0pt plus 1pt}

%==============================================================================
%  ENTORNOS DE LA OBRA
%==============================================================================

% --- Modelo de escrito judicial ------------------------------------------
\newtcolorbox{escrito}[1][]{%
  breakable,
  enhanced,
  colback=ColorFondo,
  colframe=ColorBorde,
  boxrule=0.5pt,
  arc=1pt,
  left=10pt,right=10pt,top=9pt,bottom=9pt,
  fonttitle=\sffamily\small\bfseries\color{ColorTinta},
  coltitle=ColorTinta,
  colbacktitle=ColorFondo,
  title={#1},
  before skip=10pt, after skip=10pt
}

% --- Nota de práctica forense --------------------------------------------
\newtcolorbox{practica}[1][Nota de práctica]{%
  breakable,
  enhanced,
  colback=white,
  colframe=ColorAcento,
  boxrule=0pt,
  leftrule=2.2pt,
  arc=0pt, outer arc=0pt,
  left=10pt,right=6pt,top=7pt,bottom=7pt,
  fontupper=\small,
  title={\sffamily\small\bfseries\color{ColorAcento}#1},
  colbacktitle=white, coltitle=ColorAcento,
  before skip=10pt, after skip=10pt
}

% --- Advertencia / trampa procesal ---------------------------------------
\newtcolorbox{alerta}[1][Advertencia]{%
  breakable, enhanced,
  colback=ColorFondo, colframe=ColorAcento,
  boxrule=0.7pt, arc=1pt,
  left=10pt,right=8pt,top=7pt,bottom=7pt,
  fontupper=\small,
  title={\sffamily\small\bfseries #1},
  before skip=10pt, after skip=10pt
}

% --- Cita normativa en bloque --------------------------------------------
\newenvironment{norma}[1]%
  {\begin{quote}\small\noindent\textbf{\sffamily #1}\par\vskip3pt\itshape}%
  {\end{quote}}

%==============================================================================
%  DIAGRAMAS DE FLUJO (TikZ)
%==============================================================================
\tikzset{
  nodo/.style={
    rectangle, draw=ColorBorde, fill=ColorFondo, line width=0.6pt,
    rounded corners=2pt, align=center, inner sep=6pt,
    font=\sffamily\scriptsize, text width=52mm, minimum height=9mm},
  hito/.style={nodo, fill=ColorTinta!12, draw=ColorTinta, text=ColorTinta,
    font=\sffamily\scriptsize\bfseries},
  desenlace/.style={nodo, text width=44mm},
  favorable/.style={desenlace, fill=ColorTinta!10, draw=ColorTinta},
  adverso/.style={desenlace, fill=ColorAcento!10, draw=ColorAcento},
  bifurcacion/.style={
    diamond, aspect=2.4, draw=ColorBorde, fill=white, line width=0.6pt,
    align=center, inner sep=2pt, font=\sffamily\scriptsize, text width=32mm},
  flecha/.style={draw=ColorGris, line width=0.7pt, -{Latex[length=2.2mm,width=1.6mm]}},
  rotulo/.style={font=\sffamily\scriptsize\itshape, text=ColorGris, inner sep=2pt},
}

% \begin{flujograma}{Título}{etiqueta} ... \end{flujograma}
\newenvironment{flujograma}[2]{%
  \def\FlujoTitulo{#1}\def\FlujoEtiqueta{#2}%
  \begin{figure}[htbp]\centering
  \begin{tikzpicture}[node distance=7mm and 12mm]%
}{%
  \end{tikzpicture}
  \caption{\FlujoTitulo}\label{fig:\FlujoEtiqueta}
  \end{figure}%
}
